\section{The two-wall criterion}
\label{sec:two-wall}

Let \(Y\) be an intermediate smooth projective variety, \(D\subset Y\) the
union of the two incident wall boundaries, and \(p\in U:=Y\setminus D\).
Suppose a common analytic receiver contains the two sectorial realizations
under comparison, and let \(T\) be their transition.  Write
\(\mathcal S_D(Y)\) for a span of Gamma sections represented by classes
supported on \(D\).  The point row annihilates this span:
\[
 \rrow_{Y,p}(v)=0\qquad(v\in\mathcal S_D(Y)).                  \tag{7.1}
\]

\begin{hypothesis}[Rank-zero-target condition]
\label{hyp:rank-zero-target}
On the packet under consideration, the relative transition admits a
canonical ray-ordered factorization
\[
 T=\prod_{a=1}^m(1+v_a\otimes\lambda_a)                       \tag{7.2}
\]
whose complete residue targets \(v_a\) belong to \(\mathcal S_D(Y)\).
\end{hypothesis}

The word \emph{complete} is essential.  Refining a residue block into
coordinate terms can create nonzero ranks which cancel only after the full
block is assembled.

\begin{theorem}[One-row assembly]
\label{thm:rank-zero-target}
Under Hypothesis~\ref{hyp:rank-zero-target},
\[
 \rrow_{Y,p}T=\rrow_{Y,p}.                                    \tag{7.3}
\]
Consequently, any finite factorization whose ordinary-flop steps satisfy
Theorem~\ref{thm:ordinary-flop-point-row} and whose discrepant adjacent-wall
transitions satisfy Hypothesis~\ref{hyp:rank-zero-target} preserves the
point-row Boolean on the transported packet.
\end{theorem}

\begin{proof}
Equation (7.1) gives
\[
 \rrow_{Y,p}(1+v_a\otimes\lambda_a)=\rrow_{Y,p}
\]
for every factor in (7.2).  Multiplying the identities proves (7.3).  The
factorization statement follows by induction, using
Theorem~\ref{thm:ordinary-flop-point-row} at crepant steps.
\end{proof}

A stronger geometric condition is one-sided kernel support.  If the cone of
the transition kernel relative to the diagonal has support
\[
 \Supp(K_{\mathrm{rel}})\subset Y\times D,                    \tag{7.4}
\]
with \(D\) in the output factor, then every correction target is supported on
\(D\) and Theorem~\ref{thm:rank-zero-target} applies.  Support in
\(D\times Y\) would constrain the source and would not imply (7.3).

In the Fourier-boundary language, let \(B_{\mathrm{corner}}\) be the complete
oriented correction left by the two localizations after subtracting the
common point contribution.  Its signed punctual multiplicity is the
alternating \(\delta_{00}\) coefficient
\[
 \mu_{00}(B_{\mathrm{corner}}).
\]
Fourier transform exchanges this coefficient with the zero-section
multiplicity, hence with generic output rank.  The singular form of the
rank-zero-target condition is therefore
\[
 \mu_{00}(B_{\mathrm{corner}})=0                              \tag{7.5}
\]
after applying the point-row marking.  Proposition
\ref{prop:punctual-corner} explains why merely knowing that
\(B_{\mathrm{corner}}\) is punctual proves nothing.

\begin{remark}[What remains to be proved]
Neither (7.2), (7.4), nor (7.5) is deduced here for a general pair of
discrepant walls.  Existing one-wall theorems determine formal decompositions
and, in special toric settings, Gamma/window comparisons.  The missing
statement is directed and two-wall: it must retain which side of every
vanishing cycle is the output and must compare that orientation with the
Gamma point row.
\end{remark}
